% Appendix C: Supervisory Meeting Minutes — Handbook 8.2.10
% Good practice for all students (mandatory for MSci Cybersecurity).
\section{Supervisory Meeting Minutes}
\label{app:minutes}


\subsection*{Meeting 1 — [29/01/2026]}
\begin{description}
  \item[Attendees:] Anker Rasmussen, Martin Brain
  \item[Discussion:] How to handle decryption and key delivery between
    auction participants.  Communication patterns over Veilid nodes:
    \texttt{app\_broadcast} vs.\ \texttt{app\_message} for different message
    types.  Initial discussion of what the MPC program needs to compute
    and how results are communicated back to participants.
  \item[Actions:] Investigate Veilid messaging primitives and determine
    which is appropriate for each communication pattern.  Sketch the MPC
    program requirements.
\end{description}

\subsection*{Meeting 2 — [05/02/2026]}
\begin{description}
  \item[Attendees:] Anker Rasmussen, Martin Brain
  \item[Discussion:] Secret transfer between parties is leaky: transferring
    the decryption key reveals information about who won.  Broader metadata
    leakage concerns: when should secrets be shared with which parties, and
    how to obscure participant identities throughout the protocol.
  \item[Actions:] Research metadata leakage vectors.  Design the secret
    transfer step to minimise information revealed to non-winning parties.
\end{description}
\subsection*{Meeting 3 — [12/02/2026]}
\begin{description}
  \item[Attendees:] Anker Rasmussen, Martin Brain
  \item[Discussion:] Security model for auction results: only the winner and
    seller should receive sensitive information (decryption key and bid values).
    All other metadata exposure should be minimised.
  \item[Actions:] Design the post-MPC verification flow around this principle:
    seller learns winner identity and winning bid; all other parties learn only
    whether they won or lost.
\end{description}
\subsection*{Meeting 4 — [19/02/2026]}
\begin{description}
  \item[Attendees:] Anker Rasmussen, Martin Brain
  \item[Discussion:] Arrived with semi-honest (Shamir) as the default
    protocol choice assuming honest majority.  Discussed whether to change
    the MPC auction algorithm, and what exactly needs to be computed inside
    the MPC program.  Key question: how to handle bid commitment and how to
    prove the winner actually bid what they claim.
  \item[Actions:] Define the commitment scheme (hash-based) and design a
    post-MPC verification challenge so the seller can confirm the winning
    bid matches the commitment.
\end{description}
\subsection*{Meeting 5 — [26/02/2026]}
\begin{description}
  \item[Attendees:] Anker Rasmussen, Martin Brain
  \item[Discussion:] DHT registry permissions: if all parties hold a modify
    key, any party can delete other listings.  Proposed a registry-of-registries
    pattern, but open question on how to bootstrap discovery (rebroadcast?).
    Route reuse across auctions: if a participant joins two auctions on the
    same route, this could enable deanonymisation.  MASCOT vs.\ Shamir
    selection: MP-SPDZ is somewhat opaque regarding internal behaviour.
    Party count and devnet sizing.  Initial hypothesis that MASCOT was
    network-bound on the Veilid transport layer (later found to be incorrect:
    the bottleneck was running three Tokio runtimes in a single process).
  \item[Actions:] Investigate per-listing DHT records to avoid shared-registry
    permission conflicts.  Ensure fresh routes per auction.  Separate market
    nodes into individual OS processes.
\end{description}
\subsection*{Meeting 6 — [12/03/2026]}
\begin{description}
  \item[Attendees:] Anker Rasmussen, Martin Brain
  \item[Discussion:] Martin asked to formalise the security model: protecting
    what, from whom, how, and under what assumptions.  Party assignment:
    how to deterministically assign party IDs (0, 1, 2, \ldots) given that
    party~0 (seller) is privileged.  Options discussed: sort by bid
    timestamp, sort by node ID.  Must be deterministic so all parties
    independently agree on the same assignment.
  \item[Actions:] Write up the threat model.  Implement party assignment by
    bid timestamp (seller's auto-bid is earliest, guaranteeing party~0).
\end{description}
\subsection*{Meeting 7 — [19/03/2026]}
\begin{description}
  \item[Attendees:] Anker Rasmussen, Martin Brain
  \item[Discussion:] Switched default protocol from Shamir to MASCOT, and
    from unsafe to private routing.  Veilid v0.6.0 reportedly bringing
    significant routing changes.  E2E test flakiness.  Compare-and-swap is
    not a Veilid-native primitive and can only be approximated at the
    application layer; Martin suggested using a CRDT instead.  SSL is only
    required with Shamir, not MASCOT (which uses OT-based key setup);
    system is fully reliant on Veilid's privacy layer, so if Veilid is
    compromised the MPC is also compromised.  Concurrent writes to the same
    DHT subkey caused race conditions (no concurrent primitives in Veilid).
    Veilid does not support a multi-owner model for DHT records.  The 32\,KB
    DHT record limit means a single record can realistically support at most
    256 bidders.
  \item[Actions:] Implement CRDT-based (G-Set) bid registry.  Accept Veilid
    as the sole trust boundary for transport encryption.  Design around the
    single-owner DHT constraint.
\end{description}
\subsection*{Meeting 8 — [26/03/2026]}
\begin{description}
  \item[Attendees:] Anker Rasmussen, Martin Brain
  \item[Discussion:] CRDT implementation: for true CRDTs each node needs to
    maintain its own copy, with state stored as subkeys on its own DHT
    record and a merge layer baked into the application layer.  Current
    approach was hacky with the seller centralising state.  Martin's
    feedback: local copies should converge via merge, not via a central
    authority.  Veilid's DHT still does not support multi-writer records
    cleanly.
  \item[Actions:] Refactor bid registry to use G-Set merge semantics with
    per-node owned records.  Remove seller as central state authority.
\end{description}
\subsection*{Meeting 9 — [02/04/2026]}
\begin{description}
  \item[Attendees:] Anker Rasmussen, Martin Brain
  \item[Discussion:] Dissertation writing strategy: start writing now,
    begin with the conclusion and work backwards.  Implementation priority:
    finish the CRDT refactor first, then write up results.
  \item[Actions:] Implement G-Set CRDT merge.  Begin drafting conclusion
    and results chapters.
\end{description}